\subsection{Keeping the denominator exact}

For $2\le k\le N$, define
\begin{equation}\label{eq:globalformula}
\mathcal A_{N,k}
=D_kH_k e^{Nq(\sqrt{K/N})}
(1+\sqrt{K/N})^{\epsilon-1/2}.
\end{equation}
Every factor is positive and finite, including at $k=N$. Write
\[
\mathcal E_{N,k}=\frac{R_k}{\mathcal A_{N,k}}-1.
\]
The full exponent and the parity amplitude in \eqref{eq:globalformula} give a sharper global rate than the quartically corrected approximation considered in the next section. Its largest error occurs at cubic orders.

For $c>0$, put
\begin{equation}\label{eq:globalprofiledef}
\begin{gathered}
\ell(c)=\frac23c^{3/2},\qquad
p_0(t)=t\tanh t,\qquad p_1(t)=t\coth t-1,\\
\Phi_\epsilon(c)=\frac{3(2\epsilon-1)}{4c}
\bigl(p_\epsilon(\ell(c))-\ell(c)\bigr).
\end{gathered}
\end{equation}
Both profiles extend continuously to zero with value zero.

\begin{theorem}[Sharp uniform approximation at every order]\label{thm:global}
As $N\to\infty$,
\begin{equation}\label{eq:globalerror}
\sup_{2\le k\le N}|\mathcal E_{N,k}|
=C_*N^{-1/3}+o(N^{-1/3}),
\end{equation}
where the positive absolute constant is
\begin{equation}\label{eq:globalsharpconstant}
C_*=\max_{c>0}\frac3{4c}
\left(1-\frac{2\ell(c)}{e^{2\ell(c)}-1}\right).
\end{equation}
More precisely, the full-range parity profile satisfies
\begin{equation}\label{eq:globalprofile}
\sup_{2\le k\le N}
\left|N^{1/3}\mathcal E_{N,k}
-\Phi_{k\bmod2}(k/N^{1/3})\right|\longrightarrow0.
\end{equation}
Uniformly for $\sqrt N\le k\le N$, the error is
$O(K^{-1}+e^{-\lambda/8})$, and hence $O(N^{-1/2})$.

The maximum in \eqref{eq:globalsharpconstant} is attained at a unique $c_*>0$. If $z_*$ is the unique positive root of
\begin{equation}\label{eq:globalmaxequation}
2e^{2z}-(3z^2-z+4)e^z-z+2=0,
\end{equation}
then $c_*=(3z_*/4)^{2/3}$. Numerically,
\[
c_*\approx1.16540187318,\qquad C_*\approx0.395496270091.
\]
Odd orders nearest $c_*N^{1/3}$ attain the limiting supremum, with negative signed error. In particular, the exponent $1/3$ in \eqref{eq:globalerror} is optimal for \eqref{eq:globalformula}.
\end{theorem}
\begin{proof}
Appendix~\ref{sec:globalsharpproof} proves the compact critical expansion \eqref{eq:globalcompact}, retains both real saddle contributions, and controls both tails of $k/N^{1/3}$. These estimates give \eqref{eq:globalprofile}. The odd absolute profile strictly exceeds the even profile, and the appendix proves uniqueness of its maximizer, yielding \eqref{eq:globalerror} and \eqref{eq:globalsharpconstant}.

For the upper-range bound, divide \eqref{eq:normalizedsaddle} by the exact $H_k$. The residual factor is $(1+(-1)^\epsilon e^{-2\lambda})^{-1}$, which differs from one exponentially little for $k\ge\sqrt N$. The factor $D_k$ cancels exactly. This also explains why no approximation to $(N)_k$ is needed at $k=N$.
\end{proof}

At cubic orders, the four-block term and the common quadratic availability loss cancel between the moment and \eqref{eq:globalformula}. The remaining error records the difference between the parity-conditioned triple count and its mean parameter. Its signs are opposite in the two parities. The uniform $N^{-1/4}$ rate in the next section concerns the different approximation $e^{k^2/(4N)}H_k$.

At the top order, Lemma~\ref{lem:saddle} can be used with degree $K=N$ and offset zero, without a parity restriction on that lemma. Here $r=1/2$ and $b=4N$, giving
\[
[z^N]e^{Nf(z)}=\frac{e^N}{\sqrt{8\pi N}}(1+O(N^{-1})),\qquad
a_N(N)=\frac{1+O(N^{-1})}{2N!}.
\]
This endpoint also follows from the classical truncated-exponential identity for $L_N^{(-N-1)}$. It is a consistency check on the full-range formula, rather than a separate priority claim.
